\section*{الفصل \ref{ch:b1:structure-spectroscopy} --- التحليل الطيفي لتعيين البنية}

\begin{solution}{exo:b1:structure-spectroscopy:1}
$\varepsilon = A/(\ell c) = 0.64/(1.00 \times \num{2.0e-5}) =
\qty{3.2e4}{L/(mol.cm)}$. وبثلاثة أضعاف التركيز: $A = 1.92$، إن
بقي القانون صالحًا عند هذه \omterm{def:b1:structure-spectroscopy:absorbance}{الامتصاصية}.
\end{solution}

\begin{solution}{exo:b1:structure-spectroscopy:2}
\ce{C6H6}: $D = (12 + 2 - 6)/2 = 4$، البنزين (حلقة وثلاث روابط $\pi$).
\ce{C4H8O2}: 1، إيثانوات الإيثيل أو حمض البوتانويك. \ce{C3H7NO}: $(6 + 2 + 1
- 7)/2 = 1$، البروبانأميد. \ce{C2H3Cl}: 1، الكلوروإيثين. \ce{C10H14O}:
$(20 + 2 - 14)/2 = 4$: في الكارفون حلقة واحدة، ورابطتان C=C ورابطة C=O واحدة.
\end{solution}

\begin{solution}{exo:b1:structure-spectroscopy:3}
البروبانون: إشارة واحدة (6H). حمض الإيثانويك: إشارتان (3H، 1H). 2-ميثيل البروبان:
إشارتان (9H، 1H). 1,2-ثنائي كلورو الإيثان: إشارة واحدة (4H).
\end{solution}

\begin{solution}{exo:b1:structure-spectroscopy:4}
$\delta = 1460/400 = \qty{3.65}{ppm}$؛ وعند \qty{90}{MHz} تقع الإشارة
على بعد $3.65 \times 90 = \qty{329}{Hz}$ من المرجع. وتبقى \omterm{def:b1:structure-spectroscopy:coupling}{ثوابت الاقتران}
ثابتة بالهرتز بينما تكبر فروق الانزياح مع المجال:
\omterm{def:b1:structure-spectroscopy:coupling}{فمتعددات الخطوط} المتراكبة في المجال المنخفض تنفصل في المجال العالي.
\end{solution}

\begin{solution}{exo:b1:structure-spectroscopy:5}
\ce{CH3}: ثلاثية (3H)؛ \ce{CH2} الوسطى: خمسة جيران، سداسية
1\,:\,5\,:\,10\,:\,10\,:\,5\,:\,1 (2H)؛ \ce{CH2Br}: ثلاثية (2H)، وهي الأشد
إزالةً للحجب.
\end{solution}

\begin{solution}{exo:b1:structure-spectroscopy:6}
الإيثانول: O--H (3666) ولا C=O؛ البروبانون: C=O (1738) ولا O--H؛
حمض الإيثانويك: كلتاهما، O--H (3582) وC=O (1788).
\end{solution}

\begin{solution}{exo:b1:structure-spectroscopy:7}
إيثانوات الإيثيل: رباعية (2H) قرب 4.1، وأحادية (3H) قرب 2.0، وثلاثية
(3H) قرب 1.3 جزءًا من المليون. بروبانوات الميثيل: أحادية (3H) لمجموعة \ce{OCH3}، أزال الأكسجين
الحجب عنها (قرب 3.7 جزءًا من المليون)، ورباعية (2H) لمجموعة \ce{CH2C=O} (قرب 2.3)، وثلاثية
(3H) قرب 1.1. فالرباعية هي الإشارة الأشد إزالةً للحجب في الأول
لا في الثاني: والرباعية قرب 4 أجزاء من المليون تعني \ce{O-CH2CH3}.
\end{solution}

\begin{solution}{exo:b1:structure-spectroscopy:8}
خطوط تفصل بينها \qty{0.080}{ppm}، أي $0.080 \times 90 = \qty{7.2}{Hz}$:
\omterm{def:b1:structure-spectroscopy:coupling}{ثابت الاقتران}. وللثلاثية التباعد نفسه: فمجموعتا \ce{CH2}
و\ce{CH3} مقترنتان إحداهما بالأخرى.
\end{solution}

\begin{solution}{exo:b1:structure-spectroscopy:9}
حزمة C=O وإشارة وحيدة: البروبانون، \ce{CH3COCH3}. أما المتماكب
الألدهيدي، البروبانال \ce{CH3CH2CHO}، فيُظهر حزمة اهتزاز امتطاط C--H لمجموعة
\ce{CHO} وبروتونها الشديد إزالة الحجب.
\end{solution}

\begin{solution}{exo:b1:structure-spectroscopy:10}
تشطر المجموعة الأولى الخط إلى $n_1 + 1$ خطًا؛ وتشطر المجموعة الثانية كلًّا منها
إلى $n_2 + 1$: أي $(n_1 + 1)(n_2 + 1)$ خطًا. وإذا كان $J_1 =
J_2$، فإن موضع الخط لا يتعلق إلا بالمجموع $k_1 + k_2$، الذي يأخذ $n_1 + n_2 +
1$ قيمة: وهذه قاعدة $n + 1$ مع $n = n_1 + n_2$. \ce{CH2} الوسطى في
1-برومو البروبان: $4 \times 3 = 12$ خطًا عمومًا، وستة (سداسية) حين
يتساوى الثابتان.
\end{solution}

\begin{solution}{exo:b1:structure-spectroscopy:11}
$D = 0$؛ كحول (حزمة O--H، ولا C=O). الثنائية (6H): مجموعتا \ce{CH3}
متكافئتان على CH واحدة؛ \omterm{def:b1:structure-spectroscopy:coupling}{ومتعدد الخطوط} (1H): تلك المجموعة CH؛ والثنائية (2H): مجموعة \ce{CH2}
مجاورة لمجموعة CH وتحمل OH؛ والأحادية العريضة: OH.
2-ميثيل البروبان-1-أول، \ce{(CH3)2CH-CH2-OH}.
\end{solution}

\begin{solution}{exo:b1:structure-spectroscopy:12}
كل نوع يمتص باستقلال: $A = A_1 + A_2 = \ell(\varepsilon_1c_1
+ \varepsilon_2c_2)$. وعند طولين موجيين، مع معرفة المعاملات الأربعة،
تعطي معادلتان خطيتان $c_1$ و$c_2$.
\end{solution}

\begin{solution}{pb:b1:structure-spectroscopy:1}
\textbf{1.} $0.5690 \times 12.0/44.0 = \qty{0.1552}{g}$.
\textbf{2.} $0.2328 \times 2.0/18.0 = \qty{0.02587}{g}$.
\textbf{3.} $0.2500 - 0.1552 - 0.0259 = \qty{0.0689}{g}$.
\textbf{4.} C \qty{62.1}{\percent}، H \qty{10.3}{\percent}، O
\qty{27.6}{\percent}.
\textbf{5.} C $0.1552/12.0 = 0.01293$، H $0.02587$، O $0.0689/16.0 =
0.00431$ مول: النسبة $3.00 : 6.00 : 1$، \ce{C3H6O}.
\textbf{6.} $M = \rho RT/p = 3740 \times 8.314 \times 373.15/(1.000 \times 10^5) =
\qty{116}{g/mol}$ (الكثافة بالوحدة \unit{g/m^3}).
\textbf{7.} $116/58 = 2$: \ce{C6H12O2}، $D = (12 + 2 - 12)/2 = 1$.
\textbf{8.} اهتزاز امتطاط C=O.
\textbf{9.} أي O--H: فلا حمض ولا كحول.
\textbf{10.} اهتزاز امتطاط C--O، قوي: رابطة C--O في إستر.
\textbf{11.} اهتزازات امتطاط C--H لمجموعات \ce{CH2} و\ce{CH3}.
\textbf{12.} رابطة $\pi$ واحدة، تستهلكها C=O؛ ومع C--O قوية وغياب
O--H، فهو إستر. فالحمض يلزمه O--H؛ وثنائي الأول لا C=O فيه. والكيتون
الذي يحمل مجموعة إيثر تكون حزمة C=O فيه قرب \qty{1738}{cm^{-1}}
للبروبانون، بينما تقع القيمة المرصودة \qty{1754}{cm^{-1}} قريبًا من
\qty{1764}{cm^{-1}} لإيثانوات الإيثيل؛ ويؤكد الرنين المغناطيسي النووي في الجزء الثالث
الإستر.
\textbf{13.} خمس؛ $2 + 2 + 2 + 3 + 3 = 12$، أي الهيدروجينات الاثنا عشر.
\textbf{14.} مجموعة \ce{CH2} لها ثلاثة جيران (مجموعة \ce{CH3})، مرتبطة
بالأكسجين (أُزيل الحجب عنها): \ce{-O-CH2-CH3}.
\textbf{15.} مجموعة \ce{CH3} في مجموعة الإيثيل تلك، ولها جاران؛ والإشارتان
مقترنتان إحداهما بالأخرى، ومن ثم قيمة $J$ واحدة.
\textbf{16.} مجموعة \ce{CH2} لها جاران، مجاورة للرابطة C=O.
\textbf{17.} مجموعة \ce{CH2} لها خمسة جيران: بين مجموعة \ce{CH2} ومجموعة
\ce{CH3}.
\textbf{18.} مجموعة \ce{CH3} لها جاران، بعيدة عن المجموعة الوظيفية.
\textbf{19.} \ce{CH3CH2O-} و\ce{-CH2CH2CH3}.
\textbf{20.} بروتونات 4.13 جزءًا من المليون على الكربون المرتبط بالأكسجين،
وبروتونات 2.28 جزءًا من المليون مجاورة للرابطة C=O: وكلاهما ساحب للإلكترونات.
\textbf{21.} \ce{CH3CH2-O-CO-CH2CH2CH3} أو \ce{CH3CH2-CO-O-CH2CH2CH3}. ولا
يضع على الأكسجين مجموعة \ce{CH2} ذات الجيران الثلاثة إلا الأول
(الرباعية عند 4.13 جزءًا من المليون).
\textbf{22.} بروبانوات البروبيل هو الثاني: فتكون إشارة \ce{OCH2} فيه
ثلاثية وتقع الرباعية قرب 2.3 جزءًا من المليون.
\textbf{23.} إيثانوات البوتيل، \ce{CH3CO-O-C4H9}، يُظهر أحادية (3H)
قرب 2.0 جزءًا من المليون ولا رباعية.
\textbf{24.} بوتانوات الإيثيل، \ce{CH3CH2CH2-CO-O-CH2CH3}.
\textbf{25.} الأشعة تحت الحمراء: C=O 1754، C--O 1182، C--H 2982، ولا O--H. الرنين المغناطيسي النووي: \ce{OCH2}
4.13 (رباعية)، \ce{CH2CO} 2.28 (ثلاثية)، \ce{CH2} الوسطى 1.66 (سداسية)، \ce{CH3}
الإستر 1.25 (ثلاثية)، \ce{CH3} السلسلة 0.95 (ثلاثية)، \omterm{def:b1:structure-spectroscopy:integration}{والتكاملات} 2، 2، 2، 3، 3.
\textbf{26.} \textbf{\qty{116}{g/mol}}: بوتانوات الإيثيل، \ce{C6H12O2}.
\end{solution}
